\scanpage{56}
\begin{center}4 - 36 -\hfill n\textsuperscript{o} 392\end{center}

Revenant à la démonstration de la prop. 19, on a, en vertu du lemme 4,
\[
\bigcap_{\underline{P}\in P(B)}R_{\underline{P}}=\bigcap_{\underline{p}\in P(A)}(R_{[A]})_{\underline{p}},
\]
et la conclusion résulte du n\textsuperscript{o} 2, th. 2.

PROPOSITION 20.- (i) \uline{Pour qu'un} $A$-\uline{module de type fini} $M$ \uline{soit pseudo-nul}, \uline{il faut et il suffit que} $M\otimes_A B$ \uline{soit un} $B$-\uline{module pseudo-nul}.

(ii) \uline{Si} $M$ \uline{est un} $A$-\uline{module de torsion de type fini}, $M\otimes_A B$ \uline{est un} $B$-\uline{module de torsion de type fini}, \uline{et l'on a} :
\[
(22)\qquad\unsure{\chi|\gamma}_B(M\otimes_A B)=i(\unsure{\chi|\gamma}_A(M)).
\]

(iii) \uline{Si} $M$ \uline{est un} $A$-\uline{module de type fini}, $M\otimes_A B$ \uline{est un} $B$-\uline{module de type fini}, \uline{et l'on a} :
\[
\begin{aligned}
(23)\qquad &c_B(M\otimes_A B)=i(c_A(M))\\
(24)\qquad &r_B(M\otimes_A B)=r_A(M).
\end{aligned}
\]

(i) Soient $\underline{P}$ un idéal premier de $B$, $\underline{p}=\underline{P}\cap A$ ; on a $(M\otimes_A B)_{\underline{P}}=M\otimes_A B_{\underline{P}}$, (chap. II, § 2, n\textsuperscript{o} 7, prop. 18), et d'autre part $M\otimes_A B_{\underline{P}}=(M\otimes_A A_{\underline{p}})\otimes_{A_{\underline{p}}}B_{\underline{P}}= {}=M_{\underline{p}}\otimes_{A_{\underline{p}}}B_{\underline{P}}$ ; la relation $M_{\underline{p}}=0$ est donc équivalente à $(M\otimes_A B)_{\underline{P}}=0$, (chap. II, § 4, n\textsuperscript{o} 4, lemme 4). Il suffit d'appliquer cette remarque à l'idéal $\underline{P}=(0)$ et aux idéaux $\underline{P}\in P(B)$ pour prouver (i), compte tenu du n\textsuperscript{o} 4, prop. 9.

Pour prouver (ii), nous utiliserons le lemme suivant :

LEMME 5.- \uline{Soient} $M_1$, $M_2$ \uline{deux} $A$-\uline{modules de type fini}, $f:M_1\longrightarrow M_2$ \uline{un homomorphisme injectif}. \uline{Alors le noyau de} $f\otimes1_B:M_1\otimes_A B\longrightarrow M_2\otimes_A B$ \uline{est pseudo-nul}.

Soit $\underline{p}$ un idéal premier de $A$ de hauteur $\le1$. On a $(M_i\otimes_A B)_{\underline{p}}=(M_i)_{\underline{p}}\otimes_{A_{\underline{p}}}B_{\underline{p}}$ $(i=1,2)$, (chap. II, § 2, n\textsuperscript{o} 7, prop. 18), et $(f\otimes1_B)_{\underline{p}}=f_{\underline{p}}\otimes1_{B_{\underline{p}}}$ ; l'hypothèse que $f$ est injectif entraîne qu'il en est de même de $f_{\underline{p}}$, (chap. II, § 1, n\textsuperscript{o} 4, th. 1) ; d'autre part, vu le choix de $\underline{p}$, $A_{\underline{p}}$ est un anneau principal et $B_{\underline{p}}$ un $A_{\underline{p}}$-module sans torsion de type fini, donc \uline{libre} ; on en conclut que $f_{\underline{p}}\otimes1_{B_{\underline{p}}}$ est

\margin{55}
